% HD-R01. Root reconstruction, 21 November 2026.
% Candidate construction credited to Severn; every map below is proved here.
\subsection{HD-R01: bounded product suppliers by an explicit cone correspondence}
\label{sec:HD-R01}

We use the following notation for a specified connected abelian cover
$D\to X$ of a smooth projective complex curve of genus $g\ge1$.
Choose an ordered list of characters $\psi_k$, $1\le k\le R$, each
of order $e_k\ge3$ dividing $12$. Repetitions are allowed. Put
$D_k=D/\ker\psi_k$, and let $s_k\ge2$ be the number of branch
points of $D_k\to X$. Write
\[
 h_k=2g-2+s_k,\quad H=\sum_kh_k,\quad
 S=\sum_k(s_k-2)=H-2Rg,\quad \nu=S/2.
\]
The data considered here have
\[
 \sum_k\dim H^{1,0}(D_k)_{\psi_k^t}=H/2
       \quad(t\in(\mathbb Z/12)^\times).
\]
In particular $H$ and $S$ are even. Denote by
$\mathfrak F_{12}(S,R)$ the explicit integral product-Fermat bound
proved in HD-N15, with the polarization equal to the sum of the
hyperplane classes: explicitly, $\mathfrak F_{12}(0,R)=1$ and
$\mathfrak F_{12}(S,R)=C_{12}(S,R)$ for $S>0$. Zero-dimensional
factors may be removed; the formula $C_{12}(S,R)$ is increasing
in $R$, so retaining them in this upper bound is valid. Set
\begin{equation}
 \mathfrak C(H,S,R)=(H!)^2\,12^{2H}\,
                         \mathfrak F_{12}(S,R).       \tag{R01.1}
\end{equation}
There is an effective integral cycle of dimension $H/2$ on
$P=\prod_kD_k^{h_k}$, of degree at most (R01.1) for the sum
$\eta$ of the degree-one point classes, whose projection to
\[
        \bigotimes_k\bigwedge^{h_k}H^1(D_k)_{\psi_k}
                                                        \tag{R01.2}
\]
is nonzero. The proof constructs the cycle through finite quotients,
vector bundles and their projective closures. It also gives its
effective Abel image, so no inverse Lefschetz correspondence is
needed for the later degree estimate.

\paragraph{Provenance.}
The special-fibre/product idea was located in B.~Severn's unrefereed
Gaussian determinant-$-3$ manuscript, Section~5 (Theorem~C and its
construction). That theorem is not used as a premise. The argument
below reconstructs the actual correspondences, proves their
nonvanishing, and adds the explicit integral degree estimate.
The published inputs are Schoen's symmetric-power construction
(1988, Lemmas~1.1--1.2 and Remark~1.4, pp.~5--7), the finite-cover
comparison of SGA~1, Expose XII, Theorem~5.1, proper base change,
and the Fermat input with its quantitative proof in HD-N15.

\paragraph{1. A concrete determinant quotient.}
For one slot omit $k$, write $\Gamma=\mathbb Z/e$, and set
\[
 N=\ker(\Gamma^h\xrightarrow{\mathrm{sum}}\Gamma),\quad
 G_0=N\rtimes\mathfrak S_h,\quad W=D'^h/G_0.
\]
There are actual finite morphisms
\[
 D'^h\longrightarrow W\xrightarrow{\beta}S^hX
              \xrightarrow{\pi}\operatorname{Pic}^hX.
\]
The residual deck group of $\beta$ is $\Gamma$. Taking finite-group
invariants commutes with rational cohomology. On the $\psi$-part,
$N$-invariance requires the same character on every curve factor;
$H^0$ and $H^2$ have trivial character. Permutation invariance on
the tensor of $H^1$ factors is alternating because of the Koszul
sign. Therefore
\[
 H^b(W,\mathbb C)_\psi=0\ (b\ne h),\qquad
 H^h(W,\mathbb C)_\psi
       =\bigwedge^hH^1(D',\mathbb C)_\psi.
                                                        \tag{R01.3}
\]
Chevalley--Weil gives dimension $h$ for the space inside the wedge,
so the last space is a line. The same assertion holds for every
conjugate character, including the inverse character.

Let $b_1,\ldots,b_s$ be the branch points on $X$, with nonzero
exponents $\alpha_j\in\mathbb Z/e$. The product relation in the
abelian deck group gives $\sum_j\alpha_j=0$. Put
$T_j=b_j+S^{h-1}X$ and let $j$ denote the complement of their union.
The $\psi$-part of $\beta_*\mathbb C$ is $j_!\mathcal L$, where
$\mathcal L$ has meridian monodromies $\psi(\alpha_j)$.
Here is the check at multiple divisors. A point lying over a divisor
containing $b_j$, with any positive multiplicity, has a coordinate
fixed by its inertia group. Applying that inertia to this coordinate
fixes the point upstairs and has residual image $\alpha_j$. A faithful
character has no invariant stalk there. Off these divisors, permutations
at coincident unramified coordinates belong to $G_0$, and the residual
cover is etale. This proves the stated sheaf equality on every stratum.

\paragraph{2. Localization at one line bundle.}
The degree $h=2g-2+s\ge2g$ makes $\pi$ a projective bundle. On its
fibre $|L|$, each $T_j$ cuts the hyperplane $|L-b_j|$. Every $s-1$
of these linear conditions are independent, since
$\deg(L-\sum_{j\in J}b_j)\ge2g-1$ for $|J|\le s-1$.
Riemann--Roch shows that the $s$ conditions are dependent precisely
at
\[
                 L_0=K_X+\sum_jb_j.
\]
Off $L_0$ their complement is
$(\mathbb C^*)^{s-1}\times\mathbb A^{g-1}$.
At least one, in fact every displayed meridian has nontrivial
local monodromy. The cellular complex of a circle with such a
rank-one coefficient system has invertible differential $1-\lambda$.
Duality and Kunneth then make the compactly supported cohomology
of this complement zero.

At $L_0$, the common centre is $\mathbb P^{g-1}=|K_X|$ after
adding the branch divisor. For $s\ge3$, projection from that centre
identifies the hyperplane complement with a vector bundle of rank
$g$ over
\[
 \Omega=\{[x_1:\cdots:x_s]:\ \sum x_j=0,\quad x_j\ne0\},
             \qquad\dim\Omega=s-2.
\]
This is the explicit linear projection: in a decomposition of the
underlying vector space into centre and quotient, its fibre over a
quotient line is $\operatorname{Hom}(\text{that line},\text{centre})$.
For $s=2$ the two hyperplanes coincide, and the complement is
$\mathbb A^g$.

Proper base change applied to $R\pi_*j_!\mathcal L$ now shows that
every cohomology sheaf is supported at $L_0$. Consequently restriction
to that fibre and extension by zero from its complement are
isomorphisms on the $\psi$-part in every degree. These are actual
restriction and extension maps of mixed Hodge structures. They give
\[
 H^h(W)_\psi\ \simeq\ H_c^h(W^\circ)_\psi,
       \qquad W^\circ=\beta^{-1}(|L_0|\setminus\textstyle\bigcup T_j).
                                                        \tag{R01.4}
\]
The boundary has zero inverse-character cohomology as well, by the
open--closed exact sequence and the same assertion in every degree.

\paragraph{3. The exact Fermat cover of the complement.}
For $s\ge3$, put $d_0=\gcd(e,\alpha_1,\ldots,\alpha_s)$ and
$e'=e/d_0$. The meridians generate $H_1(\Omega,\mathbb Z)$: a loop
can be capped in projective space and the cap made transverse to the
hyperplanes, away from their intersections. The resulting punctured
cap writes its homology class as a sum of meridians.
Thus the abelian cover of $\Omega$ with the given monodromy has
$d_0$ components. Each connected component is the quotient of
\[
 F_{e'}^{s-2,\circ}
     =\{\textstyle\sum Y_j^{e'}=0,\quad Y_j\ne0\}
\]
by the kernel of the character $(\alpha_j/d_0)_j$ of
$\mu_{e'}^s/\mu_{e'}$. The map is $x_j=Y_j^{e'}$.
Its equality with the specified component follows by these exact
meridian images; finite-cover comparison makes this an algebraic
isomorphism. Pullback to the vector bundle does not add monodromy,
because its fibres are contractible. This proves that the relevant
component of $W^\circ$ is the corresponding rank-$g$ vector bundle
over this finite quotient. Replacing the character by its inverse
changes both sides by the same inverse, so no sign convention is
discarded.

The map from $F_{12}^{s-2}$ given by coordinate powers $12/e'$
dominates $F_{e'}^{s-2}$. The pulled-back character has entries
$12\alpha_j/e$ modulo $12$, all nonzero, with sum zero.
For $s=2$, use $F_{12}^{0}$, a set of $12$ points, mapping equivariantly
onto the $e$ components of $W^\circ=\Gamma\times\mathbb A^g$.
For this zero-dimensional factor choose the faithful character
$(12/e,-12/e)$ of the $12$-point torsor; all its character lines have
Hodge type $(0,0)$. The direct equivariant map to the $e$ components
gives the required one-dimensional character cohomology in degree zero.
This choice also covers nonprimitive individual branch exponents.

On a coordinate boundary of any Fermat factor, a coordinate scalar
fixes the stratum and acts nontrivially through the displayed
character. Its eigencohomology is therefore zero. Filtering by
coordinate strata proves that compactly supported cohomology of
the all-coordinate-open part equals the projective character
cohomology. The Thom map for the rank-$g$ vector bundle then
identifies (R01.4) with that Fermat character line, twisted by $(-g)$.
All these identifications also hold for inverse characters.

For the product, the character part is the tensor product of these
lines. Its Fermat degree is $S$, and its Tate twist is $(-Rg)$.
The displayed balance of Hodge numbers makes every simultaneous
conjugate a $(\nu,\nu)$ line on $\prod_kF_{12}^{s_k-2}$.
HD-N15 supplies an effective irreducible integral cycle $Z$ of
dimension $\nu$, degree at most $\mathfrak F_{12}(S,R)$, whose
projection to this line is nonzero. An irreducible component can
be selected because the pairing is linear in a cycle. Its rational
class has nonzero projection to all conjugate lines once it has
nonzero projection to one.

\paragraph{4. Nonvanishing through the actual cone.}
Map $Z$ to the product of the finite quotients just described,
restrict to their all-coordinate-open parts, take inverse image in
one chosen product of the rank-$g$ bundles, and close this inverse
image in $Y=\prod_kW_k$. Call the resulting irreducible cycle $C_Z$.
It has dimension $\nu+Rg=H/2$. Its restriction to the open part is
the vector-bundle inverse image of the quotient image of $Z$.

Let $u$ be a nonzero class of the inverse-character line in $H^H(Y)$.
By (R01.4), its restriction to the special fibre is extension by zero
of a nonzero compactly supported class. Restriction to the chosen
component stays nonzero, because the deck group permutes the components
transitively and $u$ is an eigenvector. The Thom isomorphism identifies
this restriction with the inverse Fermat character. The projection
formula for the finite quotient and adjunction between extension by
zero and Borel--Moore restriction give
\[
       \langle[C_Z],u\rangle
           =c\,\langle[Z],\xi\rangle\ne0,
             \qquad c\in\mathbb Q^\times,             \tag{R01.5}
\]
where $\xi$ generates the inverse Fermat line. The finite degrees
account for $c$; none is zero. Poincare duality pairs the character
and its inverse perfectly, proving the last inequality. Boundary
components do not enter this equality because the relevant boundary
eigensheaves have already been proved acyclic.

Choose one component $V$ of the inverse image of $C_Z$ under
$P\to Y$, and take the sum $Z_P$ of its distinct
$\prod_kG_{0,k}$-translates, each with coefficient one. Every
component maps onto $C_Z$, and their common finite mapping degree
gives $(P\to Y)_*Z_P=c'[C_Z]$ for a positive integer $c'$.
Thus (R01.5) and the projection formula prove nonvanishing of
$Z_P$ on the pulled-back determinant line. This is an effective
integral cycle, with no denominator introduced by the singular
quotient $Y$.

\paragraph{5. The uniform degree calculation.}
The support of $Z_P$ is contained in the finite inverse image of
a multicone $T$ in $\prod_k|L_{0,k}|$. Its base is the image of
$Z$ under the coordinate-power maps of degree $12$ to the respective
$\mathbb P^{s_k-2}$. If $x$ is the sum of their hyperplane classes,
the image has degree at most $12^\nu\deg Z$, by the projection
formula; finite mapping multiplicities can only increase the
right-hand side.

Coning adds $g$ to each factor's projective dimension. To verify
the bound for an arbitrary cycle in a product, write its multidegrees
as $c_{a_1,\ldots,a_R}\ge0$, with $\sum a_k=\nu$. Its degree is
$\sum c_a\nu!/\prod a_k!$. The multicone has degree
$\sum c_a(H/2)!/\prod(a_k+g)!$.
Their termwise ratio is at most $(H/2)!/\nu!\le H!$.
Consequently
\[
                  \deg_x T\le H!12^H\deg Z.          \tag{R01.6}
\]

The full morphism
$\varphi:P\to\prod_kS^{h_k}X$ is finite flat of degree
$\prod_ke_k^{h_k}h_k!$: both varieties are smooth of the same
dimension, and the constituent curve and symmetric-power maps are
finite. For the divisor $x_k=b+S^{h_k-1}X$, its restriction to
$|L_{0,k}|$ is a hyperplane, and numerically
\[
       \varphi^*\Bigl(\sum_kx_k\Bigr)
             =\sum_ke_k\eta_k\ \ge\ \eta.
\]
The difference is nef. The orbit cycle $Z_P$ is a subcycle of the
effective flat inverse image of $T$, after ignoring repeated support
components; all its coefficients are one. Therefore projection formula
and nef monotonicity give
\[
 \deg_\eta Z_P
 \le\Bigl(\prod_ke_k^{h_k}h_k!\Bigr)\deg_xT
 \le12^HH!\,H!12^H\deg Z.
\]
This is (R01.1), with all degrees taken on actual integral cycles.

\paragraph{6. Effective Abel images for the character transfer.}
In the application every $D_k$ is a quotient of the fixed Galois
curve $C$, of genus $g_C$, and $G=\operatorname{Aut}(C/\mathbb P^1)$.
Use the actual pullback homomorphism $J(D_k)\to J(C)$ after Abel
summation in its $h_k$ curve factors. Its restricted theta degree
on $D_k$ is $\deg(C\to D_k)g(D_k)\le |G|g_C$.
The positive-Hermitian Cauchy--Schwarz identity for addition of
$h_k$ vectors gives on $P$
\[
 a^*\Theta_{J(C)^R}\le H|G|g_C\,\eta.
\]
The pullback of the chosen cohomological determinant tensor is
the tensor (R01.2), up to the nonzero integer trace factors of the
quotient maps and the factorial $h_k!$ from Abel summation. Indeed,
two $H^1_{\psi_k}$ factors on the same curve have product in
$H^2_{\psi_k^2}=0$, since $e_k\ge3$. Thus the wedge expansion
uses each of the $h_k$ curve factors exactly once and is the displayed
alternating determinant tensor. Thus $a_*Z_P$ has a nonzero evaluation on that
determinant tensor. Its dimension is $H/2$, since a lower-dimensional
image would make that evaluation zero. It is an effective integral
cycle of degree at most
\begin{equation}
      (H|G|g_C)^{H/2}\,\mathfrak C(H,S,R).           \tag{R01.7}
\end{equation}
This homological supplier is exactly what is paired with the invariant
divisor correspondence in the transfer. It retains both its integral
cycle and its stated polarization throughout.

\paragraph{Whole-proof checks.}
The root author read the entire proof, including the corrected
zero-dimensional character map, the integral orbit lift and every
degree comparison. An independent complete mathematical review by
/root/sixfolds reconstructed the localization, finite-cover comparison,
Thom pairing, finite flatness, multicone shifts and Abel map and found
no fault. HD-N15 has its separate whole-proof root check. The subgroup
application and the later divisor transfer have their own check events.
